%% GENERATED by `pipeline run results` - do not edit by hand.
\begin{table}[htbp]
\centering
\caption{O que os índices de validade de fato acompanham.}\label{tab:index-mechanism}
\footnotesize
\begin{tabular}{lll}
\toprule
relação & população & $\rho$ \\
\midrule
Silhouette vs.\ parcela de processos descartados como ruído & todos os 14 & -0,308 \\
Parcela descartada vs.\ acurácia no conjunto próprio da partição & todos os 14 & -0,204 \\
Silhouette vs.\ tamanho mediano de grupo & cob.\ total 14 & 0,427 \\
Silhouette vs.\ número de grupos & cob.\ total 14 & -0,325 \\
Número de grupos vs.\ cobertura & cob.\ total 14 & -0,991 \\
Silhouette vs.\ cobertura & cob.\ total 14 & 0,341 \\
\textbf{Silhouette vs.\ acurácia no conjunto comum} & cob.\ total 14 & \textbf{-0,392} \\
\textbf{DBCV vs.\ acurácia no conjunto comum} & cob.\ total 14 & \textbf{-0,057} \\
Silhouette vs.\ acurácia, apenas grupos pequenos (top-1 médio 0,623) & 233 grupos, 14 valores & -0,151 \\
Silhouette vs.\ acurácia, apenas grupos médios (top-1 médio 0,655) & 267 grupos, 14 valores & -0,009 \\
Silhouette vs.\ acurácia, apenas grupos grandes (top-1 médio 0,781) & 89 grupos, 14 valores & 0,074 \\
\bottomrule
\end{tabular}
\par\vspace{2pt}\footnotesize\raggedright A granularidade é a relação que resiste ao escrutínio neste tamanho de amostra. As linhas de cobertura são em grande medida uma reformulação dela (a correlação parcial dado o tamanho mediano de grupo é pequena), e a diferença entre as correlações de cobertura e de acurácia não é resolvível aqui -- seu intervalo bootstrap sobre as configurações inclui zero. As linhas de faixa carregam apenas tantos valores distintos de silhouette quantas são as partições, e não quantos são os grupos, de modo que constituem evidência muito mais fraca do que suas contagens de grupos sugerem.
\end{table}
